%% =====================================================================
%%  Getting the code : fichiers en ligne, conseils, alphabet du programmeur
%% =====================================================================
\liminaire{code}{Getting the code}{Every program of the book, ready to run}

\noindent\begin{minipage}[t]{\dimexpr\linewidth-52mm\relax}\vspace{0pt}
Every program printed in this book can be downloaded free of charge. Scan the QR code opposite, or
type the address \textbf{\otCodeAffiche} in your browser. The first page of each chapter also has
its own QR code, which takes you straight to the programs of that chapter.\par\vspace{3mm}
The files are organised like the book:
\begin{itemize}
\item \py{chapters/chapter-01/}, \py{chapter-02/}\ldots: the programs of the lessons and
  mini-projects, under the name given at the top of each code box;
\item \py{solutions/chapter-01/}\ldots: the solutions of the exercises, named
  \py{exercise-01.py}, \py{exercise-02.py}\ldots{} or after the file the exercise asks for.
\end{itemize}
\end{minipage}\hfill
\begin{minipage}[t]{46mm}\vspace{0pt}
\begin{tcolorbox}[enhanced,colback=otNavyTint,colframe=otNavyTint,arc=2mm,left=3.5mm,right=3.5mm,
    top=3.5mm,bottom=3.5mm,boxsep=0pt,before skip=0pt,after skip=0pt,halign=center]
  \otTuile[6mm]{telecharger}{otNavy}\par\vspace{1.2mm}
  {\otSemi\fontsize{7}{8.5}\selectfont\color{otNavy}\otCaps[11]{All the programs}}\par\vspace{2.4mm}
  {\color{otInk}\edef\otTmpQR{\noexpand\qrcode[height=24mm,level=M]{\otCodeBase}}\otTmpQR}\par\vspace{1.6mm}
  {\otRegular\fontsize{8}{10}\selectfont\color{otMuted}\otCodeAffiche}
\end{tcolorbox}
\end{minipage}

\begin{attention}[Type first, download later]
The files are there to help you when you are stuck, or to save time on a long program. But you learn
to program with your fingers: \textbf{type every program yourself} before you look at the file. The
small mistakes you make while typing\,---\,and fix\,---\,are part of the learning.
\end{attention}

\otIntertitre{Which tools?}
In this volume, you only need \textbf{Python} and \textbf{IDLE}, the editor installed with it
(chapter~1 explains how to install them). Every program was run with Python~\otVersionPython\ and
works with any version from 3.10 onwards. Later, when your programs grow, you will move to a
professional editor such as Visual Studio Code: volume~2 shows you how. Everything you learn in IDLE
works the same way there.

\otIntertitre{The programmer's alphabet}
Programs use a few characters that you rarely type in everyday life. Here are their names, in
English, and what they are used for in this volume. Take a moment to find each of them on your
keyboard.\par\vspace{2mm}

\begin{tcolorbox}[otBase,colback=white,colframe=otRule,boxrule=0.5pt,left=3mm,right=3mm]
\renewcommand{\arraystretch}{1.32}%
\begin{tabularx}{\linewidth}{@{}>{\centering\arraybackslash}p{12mm}>{\RaggedRight\arraybackslash}p{36mm}Y@{}}
\otTH{Sign} & \otTH{Name} & \otTH{Used for}\\\midrule
\py{( )} & parentheses & calling a function: \py{print("Hi")}; grouping a calculation\\
\py{" "} & double quotes & writing text: \py{"Hello"}\\
\py{' '} & single quotes, apostrophe & writing text too: \py{'Hello'}\\
\py{\#} & hash & starting a comment\\
\py{\_} & underscore & joining words in a name: \py{business\_card}\\
\py{=} & equals sign & giving a value to a variable (chapter~3)\\
\py{*} & asterisk, star & multiplication: \py{6 * 7}\\
\py{/} & slash & division: \py{10 / 4}\\
\py{\%} & percent sign & the remainder of a division (chapter~2)\\
\py{[ ]} & square brackets & lists, and picking one element (chapters~6 and 17)\\
\py{\{ \}} & curly braces & f-strings (chapter~4) and dictionaries (chapter~19)\\
\py{<  >} & less than, greater than & comparing values (chapter~7)\\
\py{:} & colon & ending the first line of a block (chapter~8)\\
\py{\textbackslash} & backslash & special characters in text, such as \py{\textbackslash n} (chapter~4)\\
\end{tabularx}
\end{tcolorbox}
